% =============================================================================
% COMPOSICIÓN DEL TEXTO
% =============================================================================

% La separación silábica está desactivada en packages.tex por requisito de
% formato. Este margen de emergencia reduce desbordes sin recurrir a \sloppy.
\setlength{\emergencystretch}{4em}

\setstretch{2.0}
\setlength{\parindent}{0pt}
\setlength{\parskip}{6pt}

% Formato de fecha: «Lima, mes, año».
\newdateformat{limaFecha}{%
    Lima, \monthname[\THEMONTH], \THEYEAR
}

% =============================================================================
% ENCABEZADOS Y PIES DE PÁGINA
% =============================================================================

\setlength{\headheight}{14.5pt}

\pagestyle{fancy}
\fancyhf{}
\fancyhead[R]{\thepage}
\renewcommand{\headrulewidth}{0pt}
\renewcommand{\footrulewidth}{0pt}

% Las primeras páginas de los capítulos no muestran el número de página.
\fancypagestyle{plain}{%
    \fancyhf{}
    \renewcommand{\headrulewidth}{0pt}
    \renewcommand{\footrulewidth}{0pt}
}

% Estilo de las páginas preliminares que no llevan numeración visible.
\fancypagestyle{unnumbered}{%
    \fancyhf{}
    \renewcommand{\headrulewidth}{0pt}
    \renewcommand{\footrulewidth}{0pt}
}

% =============================================================================
% RECURSOS GRÁFICOS Y NOMBRES EN ESPAÑOL
% =============================================================================

\graphicspath{{figures/}{figures/logos/}}

\addto\captionsspanish{%
    \renewcommand{\contentsname}{ÍNDICE GENERAL}
    \renewcommand{\listfigurename}{ÍNDICE DE FIGURAS}
    \renewcommand{\listtablename}{ÍNDICE DE TABLAS}
    \renewcommand{\bibname}{REFERENCIAS BIBLIOGRÁFICAS}
}

% =============================================================================
% TÍTULOS DE CAPÍTULOS Y SECCIONES
% =============================================================================

% Capítulos numerados: «Capítulo 1. Título», centrado y en 14 pt.
\titleformat{\chapter}[hang]
    {\normalfont\bfseries\centering\fontsize{14}{16.8}\selectfont}
    {Capítulo \thechapter.}
    {0.5em}
    {}
\titlespacing*{\chapter}{0pt}{0.32\textheight}{20pt}

% Capítulos sin número, ubicados en la parte superior de la página.
\titleformat{name=\chapter,numberless}[block]
    {\normalfont\bfseries\centering\fontsize{14}{16.8}\selectfont}
    {}
    {0pt}
    {}
\titlespacing*{name=\chapter,numberless}{0pt}{-50pt}{20pt}

\titleformat{\section}
    {\normalfont\bfseries\fontsize{13}{15.6}\selectfont}
    {\thesection}
    {1em}
    {}

\titleformat{\subsection}
    {\normalfont\bfseries\fontsize{12}{14.4}\selectfont}
    {\thesubsection}
    {1em}
    {}

% =============================================================================
% ÍNDICE GENERAL
% =============================================================================

\renewcommand{\cftdot}{.}
\renewcommand{\cftdotsep}{0.32}
\renewcommand{\cftchapleader}{\cftdotfill{\cftdotsep}}

\renewcommand{\cftchapfont}{\normalfont}
\renewcommand{\cftchappagefont}{\normalfont}

\renewcommand{\cftchappresnum}{Capítulo }
\renewcommand{\cftchapaftersnum}{.}
\setlength{\cftchapnumwidth}{5em}

% =============================================================================
% TÍTULOS DE LOS ÍNDICES
% =============================================================================

\renewcommand{\cfttoctitlefont}{\hfill\bfseries\fontsize{14}{16.8}\selectfont}
\renewcommand{\cftaftertoctitle}{\hfill}
\renewcommand{\cftlottitlefont}{\hfill\bfseries\fontsize{14}{16.8}\selectfont}
\renewcommand{\cftafterlottitle}{\hfill}
\renewcommand{\cftloftitlefont}{\hfill\bfseries\fontsize{14}{16.8}\selectfont}
\renewcommand{\cftafterloftitle}{\hfill}

\setlength{\cftbeforetoctitleskip}{-20pt}
\setlength{\cftbeforelottitleskip}{-20pt}
\setlength{\cftbeforeloftitleskip}{-20pt}

\setlength{\cftaftertoctitleskip}{20pt}
\setlength{\cftafterlottitleskip}{20pt}
\setlength{\cftafterloftitleskip}{20pt}

% =============================================================================
% FIGURAS, TABLAS Y ECUACIONES
% =============================================================================

\numberwithin{figure}{chapter}
\numberwithin{table}{chapter}
\numberwithin{equation}{chapter}

% Índice de ecuaciones.
\newcommand{\listequationsname}{ÍNDICE DE ECUACIONES}
\newlistof{myequations}{equ}{\listequationsname}
% Sin \par: forzarlo cortaba el párrafo justo debajo del display y añadía
% \parskip más una línea completa a doble espacio. El texto que sigue a una
% ecuación («donde ...») debe continuar el mismo párrafo.
\newcommand{\myequations}[1]{%
    \addcontentsline{equ}{myequations}{%
        \protect\numberline{\theequation}#1%
    }%
}

\renewcommand{\cftequtitlefont}{\hfill\bfseries\fontsize{14}{16.8}\selectfont}
\renewcommand{\cftafterequtitle}{\hfill}
\setlength{\cftbeforeequtitleskip}{-20pt}
\setlength{\cftafterequtitleskip}{20pt}
\renewcommand{\cftmyequationsleader}{\cftdotfill{\cftdotsep}}

% Prefijos de las entradas en los índices.
\renewcommand{\cfttabpresnum}{Tabla }
\renewcommand{\cfttabaftersnum}{.\space}
\setlength{\cfttabnumwidth}{5em}

\renewcommand{\cftfigpresnum}{Figura }
\renewcommand{\cftfigaftersnum}{.\space}
\setlength{\cftfignumwidth}{5em}

\renewcommand{\cftmyequationspresnum}{Ecuación }
\renewcommand{\cftmyequationsaftersnum}{.\space}
\setlength{\cftmyequationsnumwidth}{6em}

% Leyendas uniformes para tablas y figuras.
\captionsetup[table]{
    font={small,it},
    labelfont={small,bf},
    skip=10pt
}

\captionsetup[figure]{
    font={small,it},
    labelfont={small,bf},
    skip=10pt
}

% =============================================================================
% SECCIONES FINALES
% =============================================================================

\newcommand{\setupbackmatter}{%
    \titleformat{\chapter}[block]
        {\normalfont\bfseries\centering\fontsize{14}{16.8}\selectfont}
        {}
        {0pt}
        {}
    \titlespacing*{\chapter}{0pt}{-20pt}{20pt}

    \titleformat{name=\chapter,numberless}[block]
        {\normalfont\bfseries\centering\fontsize{14}{16.8}\selectfont}
        {}
        {0pt}
        {}
    \titlespacing*{name=\chapter,numberless}{0pt}{-20pt}{20pt}
}

% =============================================================================
% ESPACIADO VERTICAL DE LAS ECUACIONES
% =============================================================================
% El interlineado doble de setspace recalcula los espacios de las fórmulas cada
% vez que se emite \normalsize, de modo que fijarlos una sola vez no basta: hay
% que volver a aplicarlos después de cada cambio de tamaño de letra.

% >>> PARÁMETRO AJUSTABLE <<<
% Aire por encima y por debajo de cada ecuación. Súbelo para separarlas más del
% texto, bájalo para pegarlas. Referencia: 12pt es aproximadamente una línea.
\newlength{\espacioEcuaciones}
\setlength{\espacioEcuaciones}{6pt}

% Separación entre los renglones de una misma ecuación multilínea
% (entornos align, gather y split).
\newlength{\espacioEntreRenglones}
\setlength{\espacioEntreRenglones}{3pt}

\newcommand{\aplicaEspacioEcuaciones}{%
    % \dimexpr fuerza la lectura como dimensión; sin él, TeX tomaría el registro
    % como glue completo e intentaría imprimir el «plus ... minus ...» sobrante.
    \setlength{\abovedisplayskip}{\dimexpr\espacioEcuaciones\relax plus 2pt minus 1pt}%
    \setlength{\belowdisplayskip}{\dimexpr\espacioEcuaciones\relax plus 2pt minus 1pt}%
    \setlength{\abovedisplayshortskip}{\dimexpr0.5\espacioEcuaciones\relax plus 1pt minus 1pt}%
    \setlength{\belowdisplayshortskip}{\dimexpr0.6\espacioEcuaciones\relax plus 1pt minus 1pt}%
    \setlength{\jot}{\espacioEntreRenglones}%
}

\makeatletter
\g@addto@macro\normalsize{\aplicaEspacioEcuaciones}
\makeatother
\AtBeginDocument{\aplicaEspacioEcuaciones}
